% =============================================================================
\chapter{Processus de D\'ecision Markovien}
\label{ch:mdp}
% =============================================================================

Imaginez un robot qui apprend \`a marcher. \`A chaque instant, il observe sa
posture (l'\emph{\'etat}), choisit une force \`a appliquer \`a ses articulations
(l'\emph{action}), et re\c{c}oit un signal --- est-il rest\'e debout ou
est-il tomb\'e~? (la \emph{r\'ecompense}). Son objectif~: trouver une strat\'egie
d'actions qui maximise les r\'ecompenses accumul\'ees au fil du temps. Ce
sch\'ema --- \'etat, action, r\'ecompense, nouvel \'etat --- est universel~: il
d\'ecrit aussi bien un joueur d'\'echecs, un v\'ehicule autonome qu'un
algorithme de trading. La formalisation math\'ematique de ce sch\'ema porte un
nom~: le \emph{processus de d\'ecision markovien} (MDP), formalis\'e par
Richard Bellman dans les ann\'ees 1950.

\begin{intuition}
Un processus de d\'ecision markovien (MDP) est le mod\`ele math\'ematique central
de l'apprentissage par renforcement. Il formalise la mani\`ere dont un agent
interagit avec un environnement en prenant des d\'ecisions s\'equentielles pour
maximiser une r\'ecompense cumul\'ee \`a long terme.
\end{intuition}

% -----------------------------------------------------------------------------
\section{Introduction}
% -----------------------------------------------------------------------------

L'apprentissage par renforcement repose sur une id\'ee simple~: un agent apprend
\`a agir en interagissant avec un environnement. \`A chaque pas de temps $t$,
l'agent observe un \'etat $s_t \in \mathcal{S}$, choisit une action
$a_t \in \mathcal{A}$, re\c{c}oit une r\'ecompense $r_{t+1} \in \R$ et
transite vers un nouvel \'etat $s_{t+1}$. Le formalisme math\'ematique qui
capture cette interaction est le \emph{processus de d\'ecision markovien}.

% -----------------------------------------------------------------------------
\section{D\'efinition formelle}
% -----------------------------------------------------------------------------

\begin{definition}[Processus de d\'ecision markovien]
Un MDP est un quintuplet $\mathcal{M} = (\mathcal{S}, \mathcal{A}, P, R, \gamma)$ o\`u~:
\begin{itemize}
    \item $\mathcal{S}$ est un ensemble fini ou d\'enombrable d'\textbf{\'etats},
    \item $\mathcal{A}$ est un ensemble fini ou d\'enombrable d'\textbf{actions},
    \item $P : \mathcal{S} \times \mathcal{A} \times \mathcal{S} \to [0,1]$ est la
    \textbf{fonction de transition}~: $P(s' \mid s, a) = \PP(S_{t+1} = s' \mid S_t = s, A_t = a)$,
    \item $R : \mathcal{S} \times \mathcal{A} \times \mathcal{S} \to \R$ est la
    \textbf{fonction de r\'ecompense}~: $R(s, a, s')$ est la r\'ecompense re\c{c}ue
    lors de la transition de $s$ \`a $s'$ sous l'action $a$,
    \item $\gamma \in [0, 1)$ est le \textbf{facteur d'actualisation}.
\end{itemize}
\end{definition}

\begin{formules}[title=\textbf{Propri\'et\'e de Markov}]
La propri\'et\'e de Markov stipule que l'\'etat futur ne d\'epend que de l'\'etat
pr\'esent et de l'action courante~:
\[
    \PP(S_{t+1} = s' \mid S_0, A_0, S_1, A_1, \ldots, S_t, A_t)
    = \PP(S_{t+1} = s' \mid S_t, A_t) = P(s' \mid S_t, A_t).
\]
\end{formules}

% -----------------------------------------------------------------------------
\section{Politique}
% -----------------------------------------------------------------------------

\begin{definition}[Politique]
Une \textbf{politique} $\pi$ est une r\`egle de d\'ecision qui prescrit le comportement
de l'agent.
\begin{itemize}
    \item \textbf{Politique d\'eterministe}~: $\pi : \mathcal{S} \to \mathcal{A}$,
    o\`u $a = \pi(s)$.
    \item \textbf{Politique stochastique}~: $\pi : \mathcal{S} \times \mathcal{A} \to [0,1]$,
    o\`u $\pi(a \mid s) = \PP(A_t = a \mid S_t = s)$ avec
    $\sum_{a \in \mathcal{A}} \pi(a \mid s) = 1$ pour tout $s$.
\end{itemize}
\end{definition}

\begin{definition}[Politique stationnaire]
Une politique $\pi$ est dite \textbf{stationnaire} si elle ne d\'epend pas du temps~:
$\pi_t = \pi$ pour tout $t \geq 0$. Sauf mention contraire, toutes les politiques
consid\'er\'ees dans ce cours sont stationnaires.
\end{definition}

% -----------------------------------------------------------------------------
\section{Fonctions de valeur}
% -----------------------------------------------------------------------------

\begin{definition}[Retour actualis\'e]
Le \textbf{retour actualis\'e} \`a partir du temps $t$ est~:
\[
    G_t = \sum_{k=0}^{\infty} \gamma^k R_{t+k+1}.
\]
Le facteur $\gamma < 1$ assure la convergence de cette somme lorsque les r\'ecompenses
sont born\'ees.
\end{definition}

\begin{definition}[Fonction de valeur d'\'etat]
La \textbf{fonction de valeur d'\'etat} sous la politique $\pi$ est~:
\[
    V^\pi(s) = \E_\pi \left[ G_t \mid S_t = s \right]
    = \E_\pi \left[ \sum_{k=0}^{\infty} \gamma^k R_{t+k+1} \;\middle|\; S_t = s \right].
\]
\end{definition}

\begin{definition}[Fonction de valeur d'action]
La \textbf{fonction de valeur d'action} (ou fonction Q) sous la politique $\pi$ est~:
\[
    Q^\pi(s, a) = \E_\pi \left[ G_t \mid S_t = s, A_t = a \right].
\]
\end{definition}

\begin{remarque}
La relation entre $V^\pi$ et $Q^\pi$ est~:
\[
    V^\pi(s) = \sum_{a \in \mathcal{A}} \pi(a \mid s) \, Q^\pi(s, a).
\]
\end{remarque}

% -----------------------------------------------------------------------------
\section{\'Equations de Bellman}
% -----------------------------------------------------------------------------

\begin{theoreme}[\'Equation de Bellman pour $V^\pi$]
Pour toute politique $\pi$ et tout \'etat $s \in \mathcal{S}$~:
\[
    V^\pi(s) = \sum_{a \in \mathcal{A}} \pi(a \mid s) \sum_{s' \in \mathcal{S}}
    P(s' \mid s, a) \left[ R(s, a, s') + \gamma \, V^\pi(s') \right].
\]
\end{theoreme}

\begin{proof}
Par d\'efinition du retour actualis\'e et de la propri\'et\'e de Markov~:
\begin{align*}
    V^\pi(s) &= \E_\pi[G_t \mid S_t = s] \\
    &= \E_\pi[R_{t+1} + \gamma G_{t+1} \mid S_t = s] \\
    &= \sum_{a} \pi(a \mid s) \sum_{s'} P(s' \mid s, a)
       \left[ R(s, a, s') + \gamma \, \E_\pi[G_{t+1} \mid S_{t+1} = s'] \right] \\
    &= \sum_{a} \pi(a \mid s) \sum_{s'} P(s' \mid s, a)
       \left[ R(s, a, s') + \gamma \, V^\pi(s') \right]. \qedhere
\end{align*}
\end{proof}

\begin{theoreme}[\'Equation de Bellman pour $Q^\pi$]
Pour toute politique $\pi$, tout \'etat $s$ et toute action $a$~:
\[
    Q^\pi(s, a) = \sum_{s' \in \mathcal{S}} P(s' \mid s, a)
    \left[ R(s, a, s') + \gamma \sum_{a' \in \mathcal{A}} \pi(a' \mid s') \, Q^\pi(s', a') \right].
\]
\end{theoreme}

\begin{proof}
On d\'eveloppe de mani\`ere similaire~:
\begin{align*}
    Q^\pi(s, a) &= \E_\pi[R_{t+1} + \gamma G_{t+1} \mid S_t = s, A_t = a] \\
    &= \sum_{s'} P(s' \mid s, a) \left[ R(s, a, s') +
       \gamma \, \E_\pi[G_{t+1} \mid S_{t+1} = s'] \right] \\
    &= \sum_{s'} P(s' \mid s, a) \left[ R(s, a, s') +
       \gamma \sum_{a'} \pi(a' \mid s') Q^\pi(s', a') \right]. \qedhere
\end{align*}
\end{proof}

\begin{formules}[title=\textbf{\'Equations de Bellman --- R\'esum\'e}]
\begin{align}
    V^\pi(s) &= \sum_a \pi(a|s) \sum_{s'} P(s'|s,a)[R(s,a,s') + \gamma V^\pi(s')]
    \label{eq:bellman_v} \\
    Q^\pi(s,a) &= \sum_{s'} P(s'|s,a)[R(s,a,s') + \gamma \sum_{a'} \pi(a'|s') Q^\pi(s',a')]
    \label{eq:bellman_q}
\end{align}
\end{formules}

% -----------------------------------------------------------------------------
\section{Politique optimale et \'equations de Bellman d'optimalit\'e}
% -----------------------------------------------------------------------------

\begin{definition}[Politique optimale]
Une politique $\pi^*$ est \textbf{optimale} si pour tout \'etat $s \in \mathcal{S}$~:
\[
    V^{\pi^*}(s) \geq V^\pi(s) \quad \text{pour toute politique } \pi.
\]
On note $V^*(s) = V^{\pi^*}(s)$ et $Q^*(s,a) = Q^{\pi^*}(s,a)$.
\end{definition}

\begin{theoreme}[\'Equations de Bellman d'optimalit\'e]
Les fonctions de valeur optimales satisfont~:
\begin{align}
    V^*(s) &= \max_{a \in \mathcal{A}} \sum_{s'} P(s'|s,a)
    \left[ R(s,a,s') + \gamma V^*(s') \right] \label{eq:bellman_opt_v} \\
    Q^*(s,a) &= \sum_{s'} P(s'|s,a)
    \left[ R(s,a,s') + \gamma \max_{a'} Q^*(s',a') \right] \label{eq:bellman_opt_q}
\end{align}
\end{theoreme}

\begin{proof}
Pour l'\'equation~\eqref{eq:bellman_opt_v}, on utilise le fait que la politique
optimale choisit l'action maximisant la valeur~:
\begin{align*}
    V^*(s) &= \max_\pi V^\pi(s) \\
    &= \max_a Q^*(s,a) \\
    &= \max_a \sum_{s'} P(s'|s,a)[R(s,a,s') + \gamma V^*(s')].
\end{align*}
Pour l'\'equation~\eqref{eq:bellman_opt_q}, on injecte $V^*(s') = \max_{a'} Q^*(s',a')$
dans l'\'equation de Bellman pour $Q^*$.
\end{proof}

\begin{theoreme}[Existence d'une politique optimale d\'eterministe]
\label{thm:existence_opt}
Pour tout MDP fini, il existe au moins une politique optimale d\'eterministe
$\pi^* : \mathcal{S} \to \mathcal{A}$ d\'efinie par~:
\[
    \pi^*(s) = \argmax_{a \in \mathcal{A}} Q^*(s, a).
\]
\end{theoreme}

% -----------------------------------------------------------------------------
\section{Op\'erateur de Bellman et contraction}
% -----------------------------------------------------------------------------

\begin{definition}[Op\'erateur de Bellman]
L'\textbf{op\'erateur de Bellman} $\mathcal{T}^\pi$ pour une politique $\pi$
est d\'efini par~:
\[
    (\mathcal{T}^\pi V)(s) = \sum_a \pi(a|s) \sum_{s'} P(s'|s,a)
    [R(s,a,s') + \gamma V(s')].
\]
L'\textbf{op\'erateur de Bellman d'optimalit\'e} est~:
\[
    (\mathcal{T}^* V)(s) = \max_a \sum_{s'} P(s'|s,a)
    [R(s,a,s') + \gamma V(s')].
\]
\end{definition}

\begin{theoreme}[Contraction de l'op\'erateur de Bellman]
\label{thm:contraction}
Les op\'erateurs $\mathcal{T}^\pi$ et $\mathcal{T}^*$ sont des contractions en
norme infinie avec facteur $\gamma$~:
\[
    \norm{\mathcal{T}^\pi V - \mathcal{T}^\pi U}_\infty
    \leq \gamma \norm{V - U}_\infty.
\]
Par le th\'eor\`eme du point fixe de Banach, chacun poss\`ede un unique point fixe~:
$V^\pi$ pour $\mathcal{T}^\pi$ et $V^*$ pour $\mathcal{T}^*$.
\end{theoreme}

\begin{proof}
Pour tout $s \in \mathcal{S}$~:
\begin{align*}
    \abs{(\mathcal{T}^* V)(s) - (\mathcal{T}^* U)(s)}
    &= \abs{\max_a \sum_{s'} P(s'|s,a)[R + \gamma V(s')]
       - \max_a \sum_{s'} P(s'|s,a)[R + \gamma U(s')]} \\
    &\leq \max_a \sum_{s'} P(s'|s,a) \gamma \abs{V(s') - U(s')} \\
    &\leq \gamma \norm{V - U}_\infty.
\end{align*}
Le passage de la deuxi\`eme \`a la troisi\`eme ligne utilise l'in\'egalit\'e
$\abs{\max_a f(a) - \max_a g(a)} \leq \max_a \abs{f(a) - g(a)}$.
\end{proof}

% -----------------------------------------------------------------------------
\section{Fonction avantage}
% -----------------------------------------------------------------------------

\begin{definition}[Fonction avantage]
La \textbf{fonction avantage} sous la politique $\pi$ est~:
\[
    A^\pi(s, a) = Q^\pi(s, a) - V^\pi(s).
\]
Elle mesure l'avantage relatif de prendre l'action $a$ dans l'\'etat $s$
par rapport au comportement moyen sous $\pi$.
\end{definition}

\begin{proposition}
Pour toute politique $\pi$ et tout \'etat $s$~:
\[
    \sum_{a \in \mathcal{A}} \pi(a \mid s) \, A^\pi(s, a) = 0.
\]
\end{proposition}

\begin{proof}
$\sum_a \pi(a|s) A^\pi(s,a) = \sum_a \pi(a|s)[Q^\pi(s,a) - V^\pi(s)]
= V^\pi(s) - V^\pi(s) = 0$.
\end{proof}

% -----------------------------------------------------------------------------
\section{Exemples de MDP}
% -----------------------------------------------------------------------------

\begin{exemple}[Gridworld]
Consid\'erons une grille $4 \times 4$ o\`u un agent se d\'eplace dans les quatre
directions cardinales. L'\'etat terminal est la case $(4,4)$ avec r\'ecompense $+1$.
Chaque d\'eplacement co\^ute $-0.01$. Les mouvements contre les murs laissent
l'agent sur place. Ce probl\`eme se formalise comme un MDP avec~:
\begin{itemize}
    \item $\mathcal{S} = \{(i,j) : 1 \leq i,j \leq 4\}$, $|\mathcal{S}| = 16$,
    \item $\mathcal{A} = \{\text{haut, bas, gauche, droite}\}$,
    \item Transitions d\'eterministes (sauf aux bords),
    \item $\gamma = 0.99$.
\end{itemize}
\end{exemple}

\begin{exemple}[Bandit \`a $K$ bras]
Un bandit \`a $K$ bras est un MDP d\'eg\'en\'er\'e avec $|\mathcal{S}| = 1$
(un seul \'etat). L'agent choisit \`a chaque pas un bras $a \in \{1, \ldots, K\}$
et re\c{c}oit une r\'ecompense al\'eatoire $R_a \sim \mathcal{D}_a$. Ce cas
particulier permet d'\'etudier le dilemme exploration-exploitation de mani\`ere
isol\'ee.
\end{exemple}

% -----------------------------------------------------------------------------
\section{Impl\'ementation en Python}
% -----------------------------------------------------------------------------

\begin{codeblock}[title=\textbf{D\'efinition d'un MDP simple avec Gymnasium}]
\begin{minted}{python}
import numpy as np
import gymnasium as gym
from gymnasium import spaces

class GridWorldEnv(gym.Env):
    """Environnement Gridworld 4x4 personnalise."""
    metadata = {"render_modes": ["human"]}

    def __init__(self, size=4, gamma=0.99):
        super().__init__()
        self.size = size
        self.gamma = gamma
        self.observation_space = spaces.Discrete(size * size)
        self.action_space = spaces.Discrete(4)  # haut, bas, gauche, droite
        self.goal = (size - 1, size - 1)
        self._action_to_dir = {
            0: np.array([-1, 0]),  # haut
            1: np.array([1, 0]),   # bas
            2: np.array([0, -1]), # gauche
            3: np.array([0, 1]),  # droite
        }
        self.reset()

    def _pos_to_state(self, pos):
        return pos[0] * self.size + pos[1]

    def reset(self, seed=None, options=None):
        super().reset(seed=seed)
        self.agent_pos = np.array([0, 0])
        return self._pos_to_state(self.agent_pos), {}

    def step(self, action):
        direction = self._action_to_dir[action]
        new_pos = np.clip(self.agent_pos + direction, 0, self.size - 1)
        self.agent_pos = new_pos

        terminated = tuple(self.agent_pos) == self.goal
        reward = 1.0 if terminated else -0.01
        return self._pos_to_state(self.agent_pos), reward, terminated, False, {}
\end{minted}
\end{codeblock}

\begin{codeblock}[title=\textbf{Construction de la matrice de transition}]
\begin{minted}{python}
def build_transition_matrix(env):
    """Construit P[s, a, s'] et R[s, a, s'] pour un MDP fini."""
    nS = env.observation_space.n
    nA = env.action_space.n
    P = np.zeros((nS, nA, nS))
    R = np.zeros((nS, nA, nS))

    for s in range(nS):
        for a in range(nA):
            # Simuler la transition depuis l'etat s avec l'action a
            env.agent_pos = np.array([s // env.size, s % env.size])
            s_next, reward, _, _, _ = env.step(a)
            P[s, a, s_next] = 1.0  # transition deterministe
            R[s, a, s_next] = reward
    return P, R

env = GridWorldEnv(size=4)
P, R = build_transition_matrix(env)
print(f"Forme de P: {P.shape}")  # (16, 4, 16)
print(f"Forme de R: {R.shape}")  # (16, 4, 16)
\end{minted}
\end{codeblock}

% -----------------------------------------------------------------------------
\section{R\'esolution alg\'ebrique directe}
% -----------------------------------------------------------------------------

\begin{proposition}[R\'esolution matricielle de l'\'equation de Bellman]
Pour un MDP fini avec $|\mathcal{S}| = n$, l'\'equation de Bellman pour $V^\pi$
peut s'\'ecrire sous forme matricielle~:
\[
    \mathbf{v}^\pi = \mathbf{r}^\pi + \gamma \, \mathbf{P}^\pi \mathbf{v}^\pi
\]
o\`u $\mathbf{P}^\pi_{s,s'} = \sum_a \pi(a|s) P(s'|s,a)$ et
$\mathbf{r}^\pi_s = \sum_a \pi(a|s) \sum_{s'} P(s'|s,a) R(s,a,s')$.
La solution est~:
\[
    \mathbf{v}^\pi = (I - \gamma \mathbf{P}^\pi)^{-1} \mathbf{r}^\pi.
\]
\end{proposition}

\begin{attention}[title=\textbf{Complexit\'e de la r\'esolution directe}]
L'inversion matricielle a une complexit\'e $O(n^3)$, ce qui la rend impraticable
pour des MDP avec un grand nombre d'\'etats. C'est pourquoi on pr\'ef\`ere les
m\'ethodes it\'eratives (Chapitre~\ref{ch:programmation_dynamique}).
\end{attention}

% -----------------------------------------------------------------------------
\section{Extensions du mod\`ele MDP}
% -----------------------------------------------------------------------------

\begin{definition}[POMDP]
Un \textbf{processus de d\'ecision markovien partiellement observable} (POMDP)
est un MDP augment\'e d'un ensemble d'observations $\mathcal{O}$ et d'une
fonction d'observation $O : \mathcal{S} \times \mathcal{A} \to \Delta(\mathcal{O})$.
L'agent n'observe pas directement l'\'etat $s_t$ mais une observation $o_t \sim O(\cdot | s_t, a_{t-1})$.
\end{definition}

\begin{definition}[MDP \`a horizon fini]
Un MDP \`a \textbf{horizon fini} $H$ est un MDP o\`u l'interaction se termine
apr\`es exactement $H$ pas de temps. Le retour est~:
\[
    G_t = \sum_{k=0}^{H-t-1} \gamma^k R_{t+k+1}.
\]
\end{definition}

\begin{definition}[MDP continu]
Lorsque $\mathcal{S} \subseteq \R^n$ et/ou $\mathcal{A} \subseteq \R^m$, on parle
de MDP \`a espace d'\'etats ou d'actions continu. Les sommes sont remplac\'ees par
des int\'egrales dans les \'equations de Bellman.
\end{definition}

% -----------------------------------------------------------------------------
\section{Exercices}
% -----------------------------------------------------------------------------

\begin{exercice}[\'Equation de Bellman]
Consid\'erons un MDP \`a 3 \'etats $\{s_1, s_2, s_3\}$ avec $\gamma = 0.9$ et
la politique uniforme $\pi(a|s) = 1/2$ pour deux actions. Les transitions
et r\'ecompenses sont~:
\begin{itemize}
    \item De $s_1$, action $a_1$~: vers $s_2$ avec $R=1$~; action $a_2$~: vers $s_3$ avec $R=0$.
    \item De $s_2$, toute action~: vers $s_1$ avec $R=2$.
    \item De $s_3$, toute action~: vers $s_3$ avec $R=0$ (\'etat absorbant).
\end{itemize}
\'Ecrivez le syst\`eme d'\'equations de Bellman pour $V^\pi$ et r\'esolvez-le.
\end{exercice}

\begin{exercice}[Op\'erateur de contraction]
Montrez que si $\gamma = 0$, l'op\'erateur de Bellman d'optimalit\'e converge
en une seule it\'eration. Quelle est l'interpr\'etation de $\gamma = 0$~?
\end{exercice}

\begin{exercice}[Impl\'ementation]
Impl\'ementez un environnement \texttt{FrozenLake} personnalis\'e avec Gymnasium.
Construisez la matrice de transition et calculez $V^\pi$ par r\'esolution
matricielle pour la politique uniforme. V\'erifiez votre r\'esultat en comparant
avec l'environnement \texttt{FrozenLake-v1} de Gymnasium.
\end{exercice}

\begin{exercice}[Fonction avantage]
Montrez que pour une politique optimale d\'eterministe $\pi^*$~:
\[
    A^{\pi^*}(s, \pi^*(s)) = 0 \quad \text{et} \quad A^{\pi^*}(s, a) \leq 0
    \quad \forall a \neq \pi^*(s).
\]
\end{exercice}
